\documentclass[11pt]{article}

% Pacotes
\RequirePackage{luatex85}
\usepackage{amsmath}
\usepackage{amssymb}
\usepackage{graphicx}
\usepackage[colorlinks=true, allcolors=blue]{hyperref}
\usepackage{amsthm}
\usepackage{framed, color,xcolor}
\usepackage{stmaryrd}
\usepackage{rotating}           % Usado para rotacionar o texto
\usepackage[all,knot,arc,import,poly]{xy}   % Pacote para desenhos gráficos
% Este pacote pode conflitar com outros pacotes gráficos como o ``pictex''
% Então é necessário usar apenas um dos pacotes conflitantes
\usepackage{tikz, tikz-cd}
\usepackage{epigraph}
\usetikzlibrary{babel}
\usetikzlibrary{positioning}
\usetikzlibrary{calc}
\usepackage[stix2]{newtxmath}
\usepackage{mathtools, nccmath}
%\usepackage[usenames,dvipsnames]{xcolor}
\usepackage{amssymb, mathrsfs}
%\usepackage[numbers]{natbib}
\usepackage{fontspec}
\usepackage{enumitem} % pra substituir o símbolo do itemize por alguma coisa legal, tipo setinhas:  \begin{itemize}[label=\ding{212}]
\usepackage{pifont}
\usepackage{stmaryrd} % for \arrownot (notação de adjacência em grafo)
\usepackage{quiver} %pacote do quiver pra desenhar diagramas mais legais

%Fontes e Línguas
\usepackage[portuguese]{babel}
%\usepackage{eufrak}
\usepackage{emoji}
\setmainfont{DarwinSerif-Regular.otf}
\newfontfamily\fakebold[FakeBold=1.75]{DarwinSerif-Regular.otf}
\newfontfamily\fakeitalic[FakeSlant=0.2]{DarwinSerif-Regular.otf}
\renewcommand{\textbf}[1]{{\fakebold#1}}
\renewcommand{\textit}[1]{{\fakeitalic#1}}
\renewcommand{\bfseries}{\fakebold}
\renewcommand{\itshape}{\fakeitalic}
\newcommand{\myref}[1]{\hyperref[#1]{\namecref{#1} \ref{#1}}}




%Cores


%Comandinhos que gosto
\usepackage{mathtools}
\newcommand{\defeq}{\vcentcolon=}
\newcommand{\eqdef}{=\vcentcolon}
\newcommand{\bigslant}[2]{{\raisebox{.2em}{$#1$}\left/\raisebox{-.2em}{$#2$}\right.}}
\newcommand{\logo}{\Rightarrow}
\newcommand{\pulinho}{\vspace{.5cm}}
\newcommand{\edge}{%
  \mathrel{-}% minus as a relation
  \joinrel\joinrel % some backup
  \mathrel{-}% minus as a relation
}
\newcommand{\notedge}{%
  \mathrel{\mkern2mu}% a small advancement
  \arrownot % a short slash
  \mathrel{\mkern-2mu}% compensate
  \edge
}

\newcommand{\Aut}{\text{Aut}}
\newcommand{\Cay}{\text{Cay}}

\newcommand{\luafacil}{
\emoji{waxing-crescent-moon}
}

\newcommand{\luamedio}{
\emoji{first-quarter-moon}
}

\newcommand{\luadificil}{
\emoji{waxing-gibbous-moon}
}

\newcommand{\luaprojeto}{
\emoji{full-moon}
}

\let\varphi\phi

\newcommand{\prabaixo}[1]{\lceil #1 \rceil}
\newcommand{\pracima}[1]{\lfloor #1 \rfloor}


%Caixinhas coloridas

\theoremstyle{definition}
\newtheorem{theorem}{Teorema}
\newenvironment{teorema}{
\definecolor{shadecolor}{HTML}{f59aa3}
\begin{shaded}
\begin{theorem}
}
{
\end{theorem}
\end{shaded}
}

\theoremstyle{definition}
\newtheorem*{theorem*}{Teorema}
\newenvironment{teorema*}{
\definecolor{shadecolor}{HTML}{f59aa3}
\begin{shaded}
\begin{theorem*}
}
{
\end{theorem*}
\end{shaded}
}

\theoremstyle{definition}
\newtheorem{prop}{Proposição}
\newenvironment{proposicao}{
\definecolor{shadecolor}{HTML}{FFDBDF}
\begin{shaded}
\begin{prop}
}
{
\end{prop}
\end{shaded}
}

\theoremstyle{definition}
\newtheorem*{prop*}{Proposição}
\newenvironment{proposicao*}{
\definecolor{shadecolor}{HTML}{FFDBDF}
\begin{shaded}
\begin{prop*}
}
{
\end{prop*}
\end{shaded}
}

\theoremstyle{definition}
\newtheorem{cor}{Corolário}
\newenvironment{corolario}{
\definecolor{shadecolor}{HTML}{f59aa3}
\begin{shaded}
\begin{cor}
}
{
\end{cor}
\end{shaded}
}

\theoremstyle{definition}
\newtheorem*{cor*}{Corolário}
\newenvironment{corolario*}{
\definecolor{shadecolor}{HTML}{f59aa3}
\begin{shaded}
\begin{cor*}
}
{
\end{cor*}
\end{shaded}
}



\theoremstyle{definition}
\newtheorem{definition}{Definição}
\newenvironment{definicao}{
\definecolor{shadecolor}{HTML}{defcfc}
\begin{shaded}
\begin{definition}
}
{
\end{definition}
\end{shaded}
}

\theoremstyle{definition}
\newtheorem*{definition*}{Definição}
\newenvironment{definicao*}{
\definecolor{shadecolor}{HTML}{defcfc}
\begin{shaded}
\begin{definition*}
}
{
\end{definition*}
\end{shaded}
}

\theoremstyle{definition}
\newtheorem{example}{Exemplo}
\newenvironment{exemplo}{
\definecolor{shadecolor}{HTML}{b9e6d3}
\begin{shaded}
\begin{example}
}
{
\end{example}
\end{shaded}
}

\theoremstyle{definition}
\newtheorem*{example*}{Exemplo}
\newenvironment{exemplo*}{
\definecolor{shadecolor}{HTML}{b9e6d3}
\begin{shaded}
\begin{example*}
}
{
\end{example*}
\end{shaded}
}

\theoremstyle{definition}
\newtheorem{veronica}{Lema}
\newenvironment{lema}{
\definecolor{shadecolor}{HTML}{FFF1F2}
\begin{shaded}
\begin{veronica}
}
{
\end{veronica}
\end{shaded}
}

\theoremstyle{definition}
\newtheorem*{veronica*}{Lema}
\newenvironment{lema*}{
\definecolor{shadecolor}{HTML}{FFF1F2}
\begin{shaded}
\begin{veronica*}
}
{
\end{veronica*}
\end{shaded}
}

\newenvironment{prova}{\paragraph{Demonstração:}}{\hfill$\square$ 
\pulinho
}

\newenvironment{aviso}{
\definecolor{shadecolor}{HTML}{fffe9a}
\begin{shaded}
}{\end{shaded}}

\theoremstyle{definition}
\newtheorem{exersisio}{Exercício}
\newenvironment{exercicio}{
\definecolor{shadecolor}{HTML}{CFB5FF}
\begin{shaded}
\begin{exersisio}
}
{
\end{exersisio}
\end{shaded}
}


\theoremstyle{definition}
\newtheorem{preguta}{Pergunta}
\newenvironment{pergunta}{
\definecolor{shadecolor}{HTML}{FFF8F0}
\begin{shaded}
\begin{preguta}
}
{
\end{preguta}
\end{shaded}
}

%Tamanho do papel
\usepackage[a4paper,top=1cm,bottom=1cm,margin=3.5cm]{geometry}



%%%%%%%%%%%%%%%%%%%%%%%%%%%%%%%%%%%%%%%%%%%%%%%


%Definições só pra esse texto
\usepackage[bb=boondox]{mathalpha}



\title{Emaranhados}
\author{Violeta Mar}
\date{}


\begin{document}

\maketitle

A teoria de emaranhados (\textit{tangles}, em inglês) nasce na teoria de conectividade de grafos finitos. Aqui, tentamos desenvolver uma versão `conjuntista' dela. A principal referência que estou usando aqui é a tese de doutorado de Jay Kneip \cite{wheretofindthem}.

Na teoria usual, a definição de sistemas de separação envolve uma ordem parcial e uma involução arbitrária. No nosso caso, iremos sempre pensar em sistemas de separação definidos por conjuntos.


\section{Sistemas, emaranhados e perfis}

A categoria que trabalhamos será a seguinte.

\begin{definicao*}[Sistemas e morfismos]
  Um conjunto $S \subset \mathcal{P}(V)$ fechado para complementar vai ser chamado de sistema. Estamos interessados em como seus elementos estão contidos entre si, e como se relacionam com a operação de complementar, logo um morfismo de sistemas será uma função $f: S \rightarrow S'$ entre dois sistemas $S \subset \mathcal{P}(V)$ e $S' \subset \mathcal{P}(V')$ que preserva inclusão e complementar. $f$ é isomorfismo de sistemas se é bijeção e sua inversa é um morfismo de sistemas e neste caso escrevemos $S \cong S'$.
\end{definicao*}

\begin{pergunta}
Essa é a definição apropriada de morfismo a se considerar aqui?
\end{pergunta}

Nosso objeto principal de estudo serão os emaranhados e os perfis de sistemas.

\begin{definicao*}[Emaranhados e perfis]Seja $V$ um conjunto, e $S \subset \mathcal{P}(V)$ um subconjunto fechado para complementar. Uma função $\tau: S \rightarrow \{0,1\}$ é dita um \textit{emaranhado} se satisfaz
  \begin{itemize}[label=\ding{212}]
  \item $\tau(A) = 1$ ou $\tau(\lnot A) = 1$ e ambos não podem acontecer ao mesmo tempo, para qualquer $A \in S$ 
  \item se $A \subset B$, então $\tau(A) \leq \tau(B)$
  \end{itemize}
Se $\tau$ satisfaz ainda a propriedade
\begin{itemize}[label=\ding{212}]
  \item se $A,B \in S$ e $A \cap B \in S$, então $\tau(A \cap B) = \tau(A)\tau(B)$
\end{itemize}
então dizemos que $\tau$ é um perfil.


Geralmente dizemos que $\tau$ escolhe um $A$ se $\tau(A) = 1$. 

Alguns nomes na literatura de tangles: a primeira condição diz que $\tau$ é uma orientação; a segunda condição diz que essa orientação é consistente; a terceira condição é dita a propriedade P.
\end{definicao*}

\begin{pergunta}
  Existe uma descrição categórica para emaranhados e perfis?
\end{pergunta}

Esta definição se aproxima muito de uma das possíveis definições de ultrafiltros para álgebras de Boole. De fato, se $S$ é uma álgebra de Boole, perfis são apenas ultrafiltros.

\begin{proposicao}[Perfis são ultrafiltros quando $S$ é de Boole]
  Suponha que $S \subset \mathcal{P}(V)$ seja fechado para união e interseção, ou seja, $S$ forma uma álgebra de Boole (ou corpo de conjuntos). Então $\tau:S \rightarrow \{0,1\}$ é um perfil se e somente se é um homomorfismo de álgebras de Boole, ou seja, um ultrafiltro.
\end{proposicao}
\begin{prova}
 (Pra escrever - lembrar que ultrafiltros são apenas homomorfismos da álgebra na álgebra $\{0,1\})$
\end{prova}

\begin{pergunta}
  Qual a relação entre os emaranhados, os perfis, os filtros e os ultrafiltros de $S$?

  Resposta simples: todo ultrafiltro de $\mathcal{P}(V)$ se restringe a um perfil, e todo perfil é um emaranhado. A pergunta aqui possivelmente mais interessante é sobre os filtros e ultrafiltros de $S$ como uma ordem parcial por si.
\end{pergunta}

Todo ultrafiltro de $\mathcal{P}(V)$ se restringe a um perfil de $S$, mas nem todo perfil de $S$ se estende para um ultrafiltro.

\begin{exemplo*}[Perfil não estendível]
  Seja $V = \{1,2,3,4\}$, $S = \{A,B,C,D,\lnot A, \lnot B, \lnot C, \lnot D\}$ onde $A = \{2,3,4\}, B = \{1,3,4\}, C = \{1,2,4\}, D = \{1,2,3\}$. A função $\tau$ que seleciona apenas os complementares $\lnot A, \lnot B, \lnot C, \lnot D$ é um perfil, porém se estendermos $S$ para a álgebra de Boole $\mathcal{P}(\{1,2,3,4\})$, $\tau$ não se estende para um ultrafiltro, já que a interseção de $\lnot A, \lnot B, \lnot C, \lnot D$ é vazia.
\end{exemplo*}

Porém, podemos consertar esse problema adicionando elementos ao $V$.

\begin{definicao*}[Principal]
  Um perfil $\tau$ de $S \subset \mathcal{P(V)}$ é dito principal se existe um $v \in V$ tal que para todo $A \in S$, $\tau(A) = 1$ se e somente se $v \in A$. 
\end{definicao*}

\begin{proposicao*}[A menos de isomorfismo perfis sempre são principais]
  Dado um $S \subset \mathcal{P}(V)$ fechado pra complementar, existe um $S' \subset \mathcal{P}(V')$ tal que $S \cong S'$ e em $S'$ todo perfil é principal. O isomorfismo preserva interseções e uniões. Podemos tomar este $S'$ e este $V'$ de forma que todo perfil é principal para um único $v \in V'$. 
\end{proposicao*}
\begin{prova}
  Para cada perfil não principal $\tau$ de $S$, defina um elemento $\alpha_{\tau}$ arbitrário e único e seja $V'$ o conjunto $V$ unido dos $\alpha_{\tau}$. Para cada $A \in S$, adicione à $A$ todo $\alpha_{\tau}$ tal que $\tau$ escolhe $A$. Ou seja, defina $A' \defeq A \cup \{\alpha_{\tau} \mid \tau(A) = 1\}$ e seja $S' = \{A' \mid A \in S\} \subset \mathcal{P}(V')$. Veja que $A \subset B$ implica em $A' \subset B'$ pela consistência de $\tau$, e é claro que $A' \subset B'$ implica em $A \subset B$. Assim, o mapa $A \mapsto A'$ é uma bijeção que preserva inclusão. O fato de $\tau$ ser uma orientação ($\tau$ escolhe um e apenas um dentre $A$ e $\lnot A$) garante que $(\lnot A)' = \lnot (A')$. Concluimos que este mapa é um isomorfismo de sistemas e logo $S \cong S'$. A propriedade P de perfil garante que $A' \cap B' = (A \cap B)'$ e portanto o isomorfismo preserva interseções e uniões quando elas existem.

  Todo perfil é principal em $S'$, pois: tome um perfil $\tau'$ de $S'$. Pelo isomorfismo, obtemos um perfil $\tau$ em $S$. Os $A'$ selecionados por $\tau'$ serão exatamente os $A$ selecionados por $\tau$, portanto cada $A'$ deve conter o elemento $\alpha_{\tau}$. Logo $\tau'$ deve ser o perfil principal em relação a $\alpha_{\tau}$.

  Para que tenhamos a parte da unicidade, defina $V''$ como sendo o quociente de $V'$ pela relação de equivalência $v \cong w$ quando $\tau_v = \tau_w$, onde estes são os perfis principais de $S$ em relação a $v,w$ respectivamente. Defina $A''$ como o quociente de $A'$ e $S'' = \{A'' \mid A' \in S'\} \subset \mathcal{P}(V'')$. Veja que $\tau_v = \tau_w$ é equivalente a dizer que $v \in A'$ se e somente se $w \in A'$ para qualquer $A' \in S'$. Assim, pode-se mostrar que $A' \subset B' \iff A'' \subset B''$ e $(\lnot A')' = \lnot (A'')$, para quaisquer $A',B' \in S'$, e portanto $S' \cong S''$. É direto ver que os perfis de $S''$ se mantém todos principais, agora principais em relação a um único elemento.
\end{prova}

Um detalhe sobre a última prova é que o único momento que precisamos da propriedade P de perfis é para garantir que o isomorfismo preserva uniões e interseções. Portanto, tirando isto, o resultado é válido para emaranhados.

Este resultado nos dá uma `forma canônica' para sistemas e seus perfis. Nesta forma, $S$ se torna uma família de subconjuntos de $V$ que separa cada um destes elementos: para cada par $\alpha,\beta \in V$, existe um $A \in S$ tal que $\alpha \in A, \beta \notin A$. Porém $S$ também satisfaz a condição de que todos os seus perfis são principais.

\begin{pergunta}
Existe uma condição `conjuntista' equivalente à condição `todos os perfis de $S$ são principais'?
\end{pergunta}

Assim como definimos a topologia de Stone no espaço de ultrafiltros de uma álgebra de Boole, definimos uma topologia no espaço de perfis.

\begin{definicao*}[Espaço de perfis]
Definimos uma topologia no conjunto de perfis $\text{Per}(S) = \{\tau: S \rightarrow \{0,1\} \text{ é perfil }\}$ como sendo a gerada pela sub-base de conjuntos da forma $U_b \defeq \{\tau \in \text{Per}(S) \mid \tau(b) = 1\}$ para cada $b \in S$.
\end{definicao*}

Como $S$ é fechado pra complementar e $\tau(b) = 1$ se e somente se $\tau(\lnot b) = 0$, temos que $U_b = \lnot U_{\lnot b}$. Assim, os geradores da topologia são fechados-abertos, e portanto $\text{Per}(S)$ é zero-dimensional. Também, conseguimos separar dois $\tau, \tau'$ distintos através do $U_b$ onde $b$ é um elemento tal que $\tau(b) \neq \tau'(b)$; ou seja, $\text{Per}(S)$ é Hausdorff. Finalmente, $\text{Per}(S) \subset \{0,1\}^S$ com a topologia induzida de subespaço. Como $\{0,1\}^S$ é compacto e 



\begin{pergunta}
  Todo $\text{Per}(S)$ é compacto?
\end{pergunta}

Se isso fosse verdade, então na verdade os espaços $\text{Per}(S)$ são apenas espaços de Stone.

\begin{pergunta}
  Se não, quem são os espaços $\text{Per}(S)$? A categoria formada por eles exibe um funtor de volta para a categoria de sistemas?
\end{pergunta}

Defina $B(S)$ como sendo a álgebra de Boole gerada por $S \subset \mathcal{P}(V)$. Seja $\text{Ult}(B(S))$ seu espaço de ultrafiltros.

\begin{proposicao}[Inclusão do espaço de ultrafiltros no espaço de perfis]
  Defina o mapa $i: \text{Ult}(B(S)) \rightarrow \text{Per}(S)$ que leva um ultrafiltro $u: B(S) \rightarrow \{0,1\}$ na sua restrição $u|_{S}: S \rightarrow \{0,1\}$. Então $i$ é um homeomorfismo sobre sua imagem.
\end{proposicao}
\begin{prova}
  Como $S$ gera $B(S)$, $i$ é injetivo. A continuidade vem do fato de que pré inversa de básico vai ser básico. Pra concluir que é homeomorfismo sobre a imagem, usamos que uma bijeção contínua de um compacto em um Hausdorff é homeomorfismo.
\end{prova}


\section{Conjuntos-Árvore}

Conjuntos-árvore vão ser sistemas de separação que capturam a estrutura de separação de uma árvore (um grafo sem ciclos).

\begin{definicao*}[Encaixados, conjunto-árvore]
  Dizemos que dois conjuntos $A,B \subset V$ estão encaixados se ou $A$ é comparável com $B$ em relação a inclusão (isto é, $A \subset B$ ou $B \subset A$) ou então $\lnot A$ é comparável com $B$ (isto é, $\lnot A \subset B$ ou $B \subset \lnot A$).

  Um $S \subset \mathcal{P}(V)$, fechado para complementar e que não contém o vazio, é dito um conjunto-árvore se quaisquer dois $A,B \in S$ (que não sejam complementares um do outro) estão encaixados.
\end{definicao*}

Lembrando que estamos sempre conceitualizando um $A \subset V$ como uma partição de $V$ em dois pedaços: $A$ e seu complementar $\lnot A$.

Vamos entender o que tem de árvore em conjuntos árvore. 

Quando falamos árvore, queremos dizer um grafo simples e sem ciclos (não necessariamente finito). Quando falamos árvore de ordem, queremos dizer um conjunto parcialmente ordenado $T$ tal que o conjunto de todos os elementos menores que $t$ é bem ordenado, pra todo $t \in T$. Vamos relaxar a definição de árvore de ordem, e definir algo que vamos chamar de \textit{arvoreta}: $T$ é uma arvoreta quando é um conjunto parcialmente ordenado tal que todos os elementos menores que $t$ formam uma cadeia (conjunto totalmente ordenado), pra todo $t \in T$.

Arvoreta na verdade é prefix order (Gustavo). Ajeitar depois.

Primeiramente, toda árvore define um conjunto árvore.

\begin{proposicao}[Conjunto árvore da árvore]
  Seja $T$ uma árvore (um grafo simples sem ciclos). Defina $S(T) = \{A \subset V(T) \mid \text{existe uma única aresta entre } A \text{ e } \lnot A\}$. Então $S$ é um conjunto árvore.
\end{proposicao}
\begin{prova}
Veja que cada aresta $e$ de $T$ define um $A \in S(T)$, tomando $A$ como uma das duas componentes conexas de $T \setminus e$, sendo a outra $\lnot A$. E, por outro lado, dado um $A \in S(T)$, se tomarmos $e$ como a aresta que liga $A$ a seu complementar, então $A$ é uma das componentes conexas de $G \setminus e$. Sabendo isso, vamos agora mostrar que $S(T)$ é um conjunto árvore. Tome dois $A,B \in S(T)$ não complementares, e seja $e,f$ suas arestas correspondentes. $f$ deve ligar dois vértices de uma das componentes conexas de $G\setminus e$, ou seja, $f$ é uma aresta do grafo induzido por $A$ ou uma aresta do grafo induzido por $\lnot A$. Digamos que $f$ é uma aresta de vértices em $A$. Então, ao retirarmos $f$, $\lnot A$ continua conexo. Logo $\lnot A$ será um subconjunto de uma das componentes conexas de $G \setminus F$, isto é, de $B$ ou $\lnot B$. Igualmente, se $f$ é uma aresta de vértices em $\lnot A$, então teriamos que $A$ é subconjunto de $B$ ou $\lnot B$. Concluimos que $S(T)$ é um conjunto árvore.
\end{prova}

Toda arvoreta define um conjunto árvore. 
Defina $\lfloor v \rfloor = \{w \in T \mid v \leq w\}$ pra um $v \in T$.

\begin{proposicao}[Conjunto árvore de uma arvoreta]
  Seja $T$ uma arvoreta. Defina $S(T) = \{A \subset T \mid A \text{ ou } A^c \text{ é da forma } \lfloor v \rfloor \text{ para algum } v \in T \}$. Então $S(T)$ é conjunto árvore.
\end{proposicao}
\begin{prova}
É uma continha simples!
\end{prova}

\begin{proposicao}[Perfis de árvores finitas são principais e recuperam a árvore]
  \label{perfis-arvore-finito}
  
  Seja $T$ uma árvore finita, e $S(T)$ seu conjunto árvore. Todos os emaranhados de $S(T)$ (e portanto todos os perfis de $S(T)$) são da forma $\tau_v$ para um $v \in V(T)$, onde $\tau_v(A) = 1 \Leftrightarrow v \in A$ para todo $A \in S$. Ou seja, $\tau_v$ é o perfil principal em relação a $v$. 

  \pulinho

  E mais: as adjacências e a métrica da árvore são recuperadas pela informação contida nos perfis. Isto é, dados dois vértices $v,w \in T$, sua distância na árvore é $n$ se e somente se seus perfis $\tau_v$ e $\tau_w$ são distintos em exatamente $n$ elementos $A \in S(T)$. No caso $n = 1$, temos que dois vértices são adjacentes se seus perfis são distintos em exatamente um elemento.
\end{proposicao}
\begin{prova}
Enumere as arestas de $T$ como $\{e_1,\dots, e_k\}$. Lembre-se da descrição na demonstração anterior de $S(T)$ via as componentes conexas de cortes da árvore. Dentre as duas componentes conexas de $T \setminus e_1$, $\tau$ seleciona uma delas, digamos $C_1$. Para quaisquer outras arestas $\{e_{i_1}, \dots, e_{i_m}\}$ na outra componente conexa, $\lnot C_1$, a consistência de $\tau$ implica que ela selecionará componentes conexas $\{C_{i_1}, \dots, C_{i_m}\}$ que vão conter $C_1$. Logo $C_1 \cap C_{i_1} \cap \dots \cap C_{i_m} = C_1$. Se $C_1$ já é um conjunto contendo apenas um vértice $v$, então esta última interseção é a interseção de todos os conjuntos selecionados por $\tau$, e portanto $\tau$ será o $\tau_v$ deste $v$. Caso contrário, devemos ter pelo menos uma aresta $e_j$ entre vértices de $C_1$. Seja $C_j$ a componente conexa escolhida por $\tau$ de $T \setminus e_j$. Olhando agora apenas ao $C_1$, ao retirarmos $e_j$ iremos ficar com duas componentes conexas, ambas não vazias: $C_1 \cap C_j$ e $C_1 \cap \lnot C_j$. Se $C_1 \cap C_j$ conter apenas um vértice, então $\tau$ é o principal em relação a este vértice. Caso contrário, ela deve conter alguma aresta, e assim poderiamos continuar este processo, encontrando a componente conexa escolhida por $\tau$ na retirada dessa aresta, e interceptando-a com $C_1$ e $C_j$. Como existem apenas finitas arestas, em algum momento esta interceção deve conter apenas um vértice, que será o $v$ tal que $\tau = \tau_v$.

Agora, se dois vértices $v,w \in T$ são adjacentes, digamos pela aresta $e_i$, então com certeza $v$ e $w$ estão em componentes conexas diferentes ao remover $e_i$, assim $\tau_v$ e $\tau_w$ diferem no corte relativo a esta aresta. E, na retirada de qualquer outra aresta $e_j$, a $e_i$ se mantém, logo os vértices estarão na mesma componente conexa e portanto $\tau_v$ e $\tau_w$ concordaram no corte relativo a $e_j$. Já se eles não forem adjacentes, deve haver um caminho entre eles contendo no mínimo duas arestas. A retirada de ambas estas arestas coloca os vértices em componentes conexas diferentes, e portanto $\tau_v$ e $\tau_w$ diferem em pelo menos dois valores. De fato, a quantidade de diferenças é exatamente o número de arestas do caminho entre os vértices, recuperando a métrica da árvore.
\end{prova}

Em uma arvoreta $T$, uma cadeia inicial $C \subset T$ é um conjunto totalmente ordenado fechado pra baixo. Uma cauda de $C$ é alguma interseção $\pracima{t} \cap C$ pra algum $t \in C$

\begin{proposicao}[Perfis de árvore de ordem correspondem a cadeias iniciais]
  Seja $\tau$ um perfil de $S(T)$, onde $T$ é uma árvore de ordem. Então existe uma cadeia inicial $C \subset T$ tal que $\tau(A) = 1$ se e somente se existe uma cauda $\pracima{t} \cap C$ contida em $A$. Cada cadeia inicial define um único perfil dessa forma.
\end{proposicao}
\begin{prova}

  Detalhe: aqui estou usando os resultados do Handbook que eu acredito que foram pensados para árvores de ordem usuais, mas que eu acredito que podem ser generalizados para arvoretas (inclusive, com as próprias demos dele, que não acredito que usam a  boa ordem dos conjuntos pra baixo).

Vou procurar uma demonstração elementar, mas por agora: extenda $S(T)$ pra $B(T)$, a tree algebra de $T$ como definida no Handbook of Boolean Algebras. Por 16.4, $\tau$ se extende pra um (único) ultrafiltro. Por 16.11, ultrafiltros correspondem a cadeias iniciais.

\end{prova}

Checar se esse resultado vale pra prefix orders - conversa com Gustavo onde vimos que não é a cadeia inicial (prefix order sendo a reta, o espaço de perfis é o conjunto de Cantor acho).

O espaço de perfis aqui será homeomorfo ao espaço de Stone de $B(T)$. Quando $T$ é um grafo, $B(T)$ é a álgebra de cortes de $T$ e logo esse espaço de Stone é homeomorfo ao espaço de direções via arestas de $T$. Logo perfis de $T$ são vértices ou raios.


% \begin{pergunta}
%   O que acontece no caso de $T$ infinita?

%   Conjectura mais ou menos intuitivamente demonstrada: os perfis de $S(T)$ são todos ou ultrafiltros principais ou ultrafiltros de extremidades de $T$.
% \end{pergunta}


% \begin{pergunta}
%   Como definir $S(T)$ para uma árvore de ordem $T$? E quem são seus perfis?

%   Conjectura mais ou menos intuitivamente demonstrada: faça do jeito mais óbvio possível, e vai dar certo, e seus perfis vão ser ultrafiltros principais ou ramos maximais.
% \end{pergunta}

Voltando ao caso finito. Temos em \ref{perfis-arvore-finito} uma descrição de como $S(T)$ contém a estrutura da árvore $T$. Se estamos de posse de seu $S(T)$, poderiamos reconstruir $T$ tomando como conjunto de vértices os perfis de $S(T)$ e ligando eles entre si quando diferem em apenas um valor. Este seria um caminho para mostrar que todo conjunto árvore finito $S$ vem de uma árvore: usando os perfis de $S$ como vértices para construir a árvore. Mas, para provar esta proposição, não vamos fazer isso. Vamos construir uma árvore explicitamente em um processo indutivo (porque essa foi a primeira ideia que eu tive).

\begin{proposicao}[Árvore do conjunto árvore]
 Todo conjunto árvore finito $S$ admite um isomorfismo (de sistemas) a um $S(T)$ para uma única árvore.
\end{proposicao}
\begin{prova}
Vamos construir uma árvore $T$ para um dado $S$ conjunto árvore finito de forma que $S \cong S(T)$. Tome alguma ordenação $S = \{A_1,\lnot A_1, A_2, \lnot A_2, \dots, A_n \lnot A_n\}$. Comece a árvore $T$ por uma árvore $T_1$ composta por dois vértices $e_1^-$ e $e_1^+$ e e uma aresta direcionada de $e_1^-$ para $e_1^+$ que nomearemos por $e_1$.

Para um $i < n$, seja $S_i = \{A_1,\lnot A_1, A_2, \lnot A_2, \dots, A_i, \lnot A_i\}$. Suponha que já construimos uma árvore $T_i$ que satisfaz

\begin{itemize}[label=\ding{212}]
  \item $T_i$ possui $i$ arestas direcionadas $\{e_1, \dots, e_i\}$
  \item $S_i \cong S(T_i) = \{\tilde{A_1},\lnot \tilde{A_1}, \tilde{A_2}, \lnot \tilde{A_2}, \dots, \tilde{A_i}, \lnot \tilde{A_i}\}$, onde cada $\tilde{A_j}$ é a imagem do $A_j$ de $S_i$
  \item $\tilde{A_j}$ é a componente conexa que contém o $e_j^+$, e portanto $\lnot \tilde{A_j}$ é a componente conexa que contém o $e_j^-$
\end{itemize}  

Já temos $T_1$ satisfazendo essas condições. Continuando o processo indutivo, vamos agora adicionar uma aresta no lugar certo de $T_i$ para construir $T_{i+1}$. Para isso, veja que o elemento $A_{i+1}$ define um emaranhado $\tau_{i+1}$ de $S_i$ da seguinte maneira.
\[ \tau_{i+1}(A_j) = 1 \text{ se } A_{i+1} \subset A_j \text{ ou } \lnot A_{i+1} \subset A_j \]
\[ \tau_{i+1}(\lnot A_j) = 1 \text{ se }  A_{i+1} \subset \lnot A_j \text{ ou } \lnot A_{i+1} \subset \lnot A_{j} \]
Como $S$ é conjunto árvore, um (e apenas um) destes dois casos deve ocorrer, dando uma função $\tau_{i+1}$ bem definida. Resta mostrar consistência. Faça essa conta.

Temos então um emaranhado $\tau_{i+1}$ de $S_i$, que naturalmente define um emaranhado $\tilde{\tau}_{i+1}$ de $S(T_i)$. Pela proposição anterior, sabemos que $\tau_{i+1}$ é principal, ou seja, existe um vértice $v \in T_i$ tal que $\tilde{\tau}_{i+1}(A) = 1 \Leftrightarrow v \in A$. Sejam $\{e_{k_1}, \dots, e_{k_m}\}$ as arestas adjacentes a $v$.``Abra'' o vértice $v$ em dois vértices $e_{i+1}^-$ e $e_{i+1}^+$ e uma aresta direcionada $e_{i+1}$ de $e_{i+1}^-$ para $e_{i+1}^+$. Para cada $k$ dentre os $k_1,\dots, k_m$, faça: 

\begin{itemize}[label=\ding{212}]
  \item se $A_{i+1} \subset A_k$, coloque o $e_k$ apontando para $e_{i+1}^-$ 
  \item se $\lnot A_{i+1} \subset A_k$, coloque o $e_k$ apontando para $e_{i+1}^+$ 
  \item se $A_{i+1} \subset \lnot A_k$, coloque o $e_k$ saindo de $e_{i+1}^-$ 
  \item se $\lnot A_{i+1} \subset \lnot A_k$, coloque o $e_k$ saindo de $e_{i+1}^+$ 
\end{itemize}

Assim temos nossa nova $T_{i+1}$, e seja $S(T_i) = \{\tilde{A_1},\lnot \tilde{A_1}, \tilde{A_2}, \lnot \tilde{A_2}, \dots, \widetilde{A_{i+1}}, \lnot \widetilde{A_{i+1}}\}$ - veja que aqui estamos cometendo um abuso de notação, repetindo os mesmos nomes $\tilde{A_j}$ que antes eram as partições da árvore anterior $S(T_i)$. Isto não é grande problema: a única diferença entre estes conjuntos é que o elemento $v$ é substituido por dois elementos $e_{i+1}^-$ e $e_{i+1}^+$. Precisamos nos convencer que $S_{i+1} \cong S(T_{i+1})$. As relações entre os novos $\tilde{A}_j$ para $j < i+1$ não mudam, pois apenas temos a troca de $v$ por $e_{i+1}^-$ e $e_{i+1}^+$ em todos os conjuntos. Basta mostrar, então, que $\widetilde{A}_{i+1}$ é encaixado com cada $\tilde{A}_j$ da mesma maneira que $A_{i+1}$ é encaixado com cada $A_j$. Tome um $j < i + 1$ e vamos fazer o caso $A_{i+1} \subset A_j$ - os outros três são análogos.

Se a aresta $e_j$ referente ao $\tilde{A}_j$ em $T_i$ era adjacente a $v$, então por construção $e_j$ aponta para $e_{i+1}^-$. Assim, o único caminho de $e_{i+1}^+$ para $e_j^+$ é a aresta $e_{i+1}$. Logo, qualquer caminho da componente conexa $\tilde{A}_{i+1}$ de $T_{i+1} \setminus e_{i+1}$ não passa pela aresta $e_k$, o que nos permite concluir que $\tilde{A}_{i+1} \subset \tilde{A}_k$.

Já se $e_j$ não for adjacente a $v$, deve existir algum caminho de $e_j^+$ até $v$. Este caminho, ao final, chega em alguma aresta $e_k$ adjacente a $v$. Gostariamos que $e_k$ fosse colocada em $e_{i+1}^-$ na construção do $T_{i+1}$, pois se isso acontece, então o único caminho de $e_{i+1}^+$ até a aresta $e_j$ passa por $e_{i+1}$ e aí conseguiriamos concluir que $\widetilde{A}_{i+1} \subset \tilde{A}_j$. Suponha, então, que $e_k$ não é colocada em $e_{i+1}^-$ na construção do $T_{i+1}$, e sim colocada em $e_{i+1}^+$. Isto só ocorre se temos que $\lnot A_{i+1}$ é subconjunto de $A_k$ ou de $\lnot A_k$. Se $\lnot A_{i+1} \subset A_k$, então $e_k$ aponta na direção de $e_{i+1}^+$. Assim, a componente conexa da ponta $e_k^+$ deve ser a que não contém a aresta $e_j$ e portanto $A_k \subset A_j$.  Mas daí teriamos $\lnot A_{i+1} \subset A_k \subset A_j$ o que contradiz $A_{i+1} \subset A_j$. Já se $\lnot A_{i+1} \subset \lnot A_k$, então $e_k$ aponta na direção oposta de $e_{i+1}^+$. Mas daí teriamos $A_k \subset \lnot A_{i+1} \subset A_j$ e também teriamos que $e_j^-$ possui um caminho até $e_k^+$ que não passa por $k$, ou seja, $e_j^- \in A_k \subset A_j$, contradição. Concluimos que $e_k$ realmente é colocada em $e_{i+1}^-$ e assim chegamos a $\widetilde{A}_{i+1} \subset \tilde{A}_j$, finalizando nossa demonstração.

Na verdade, resta provar a unicidade: se começarmos com um $S = S(T)$ para alguma árvore $T$, a árvore construida pelo algoritmo anterior é isomorfa a $T$. A demonstração disso segue facilmente se entendemos como os emaranhados de uma árvore finita recuperam a sua estrutura, como descrito na proposição anterior. De fato, suponha que $S(T_1) \cong S(T_2)$. Dado um vértice $v \in T_1$, seu emaranhado $\tau_v$ irá definir um emaranhado em $S(T_2)$, que também deverá ser principal, digamos referente ao vértice $v'$. Mapeie $\phi: v \mapsto v'$. Como também todo vértice de $T_2$ define um emaranhado que pode ser traduzido para um emaranhado de $T_1$, temos que $\phi$ é bijetora. Ela será um morfismo de grafos, isto é, $v \edge w \Leftrightarrow \phi(v) \edge \phi(w)$, pois adjacência de dois vértices acontece se e somente se seus emaranhados diferem por apenas um vértice, como descrito na proposição anterior. 

\end{prova}

\begin{proposicao}[Conjunto-árvores infinitos vem de arvoretas]
  Seja $S$ um conjunto árvore, agora possivelmente infinito. Existe uma arvoreta $T$ tal que $S \cong S(T)$.
\end{proposicao}
\begin{prova}
  Pra escrever: adaptar a demo anterior fazendo indução transfinita e usando o resultado que perfis vem de cadeias iniciais.
\end{prova}

É única? Imagino que não, por liberdade na escolha de raiz, mas talvez definam o mesmo espaço (ver pergunta abaixo). Existe uma noção de `isomorfismo' de árvore de ordem (arvoreta) a menos de escolha de raiz?

\begin{pergunta}
Tem um artigo \cite{arvoredoconjuntoarvore} que é demonstrado que conjuntos árvore infinitos se originam a partir de $T$ "tree-like space". Esses espaços são os que se originam de árvores de ordem?
\end{pergunta}

\begin{pergunta}
  No caso $S$ infinito, quem são os espaços $\text{Per}(S)$ quando $S$ é conjunto árvore?  

  Parece que vão ser os espaços de ultrafiltros de tree algebras. Que no caso acho que são espaços de cadeias iniciais de arvoretas. (topologia induzida como subespaço do espaço de caminhos?)
\end{pergunta}


\pulinho


\begin{exemplo*}
 Seja $V = \{1,2,3,4\}$, $S = \{A,B,C,D,\lnot A, \lnot B, \lnot C, \lnot D\}$ onde $A = \{2,3,4\}, B = \{1,3,4\}, C = \{1,2,4\}, D = \{1,2,3\}$. $S$ é um conjunto árvore, e se aplicarmos o algoritmo da proposição anterior, obteremos uma árvore $T$ em formato de estrela, ou seja, um vértice central com quatro arestas ligadas a ele. Temos então $5$ emaranhados nesta árvore, todos principais, relativos aos $5$ vértices de $T$. Note que o principal do vértice central não será principal em $S$: ele se traduzirá no emaranhado que escolhe sempre os complementares $\lnot A, \lnot B, \lnot C, \lnot D$, cuja interseção de todos os conjuntos escolhidos é vazia.
\end{exemplo*}

\begin{exemplo*}[Visualização]
  Pra escrever - dado um conjunto árvore $S(T)$, podemos desenhar o $V$ em formato da árvore $T$ com as separações de $S(T)$ explicitadas.
\end{exemplo*}


\section{Teorema da Árvore do Emaranhados}


Precisamos de uma condição sobre $S$ para conseguirmos enunciar e provar o teorema de Árvore de Emaranhados.

\begin{definicao*}[Sistemas sub-Boole]
  Um subconjunto $S \subset \mathcal{P}(V)$ é dito sub-Boole (ou, na terminologia do Diestel, um conjunto submodular) se é fechado por complementar e dados quaisquer dois $A,B \in S$ temos que $A \cup B \in S$ ou $A \cap B \in S$.
\end{definicao*}

O seguinte conceito será útil na demonstração do teorema.

\begin{definicao*}[Cantos, e cantos opostos]
    Dados dois conjuntos $A,B \in S$, chamamos os quatro conjuntos $A \cap B, A \cap \lnot B, \lnot A \cap B, \lnot A \cap \lnot B$ de cantos do par $A,B$. Temos dois pares de cantos que dizemos \textit{opostos}: $A \cap B$ e $\lnot A \cap \lnot B$; $\lnot A \cap B$ e $A \cap \lnot B$. 
  \end{definicao*}


Os cantos de um par de conjuntos em $S$ podem não estar presentes em $S$.
Dado um $S \subset \mathcal{P}(V)$ fechado pra complementar, dizer que dentre dois cantos opostos, pelo menos um está em $S$ é equivalente a dizer que $S$ é sub-Boole. Já dizer que todos os cantos de $S$ sempre estão em $S$ é equivalente a dizer que $S$ é uma álgebra de Boole. 

Adicionando cantos, então, sempre conseguimos estender um dado $S$ para um $\tilde{S}$ que seja um sub-Boole ou uma álgebra de Boole. Porém, um perfil $\tau$ de $S$ não necessariamente se estende para um perfil $\tilde{\tau}$ em $\tilde{S}$.

É claro, como o valor de um perfil $\tau$ em qualquer um dos quatro cantos de $A,B$ é determinado pelos valores de $\tau$ em $A$ e $B$, se existir uma extensão $\tilde{\tau}$ para $\tilde{S}$ então ela deve ser única.

Agora, a última definicação necessária antes do teorema...

\begin{definicao*}[Exibição de perfil]
  Dizemos que um $S' \subset S \subset \mathcal{P}(V)$ exibe um conjunto $\mathcal{E}$ de perfis de $S$, se para quaisquer dois $\tau, \tau' \in \mathcal{E}$, existe um $A \in S'$ tal que $\tau(A) \neq \tau'(A)$.
\end{definicao*}

\begin{teorema*}[Teorema da Árvores de Emaranhados - finito]
  Seja $V$ finito. Em qualquer $S \subset \mathcal{P}(V)$ sub-Boole existe um conjunto árvore $N \subset S$ que exibe todos os perfis de $S$.
\end{teorema*}
\begin{prova}
Vamos apresentar um algoritmo que constrói o $N$. Suponha que já construimos um conjunto árvore $N$ que diferencia um conjunto $\mathcal{O}$ de $n -1$ perfis de $S$, e ainda resta mais perfis para diferenciar. Escolha um perfil $\tau$ que $N$ não diferencia. Deve, então, existir algum outro $\tau' \in \mathcal{O}$ que concorda com $\tau$ em todos os elementos de $N$. Qualquer outro $\tau'' \in \mathcal{O}$ é diferenciado de $\tau'$ em algum elemento de $N$ por suposição, então é diferenciado de $\tau$ também. Ou seja, o único par que precisamos consertar é o $\tau$ e $\tau'$. Para isso, adicione a $N$ algum elemento $A \in S$ em que $\tau$ e $\tau'$ diferem - que, é claro, existe, pois $\tau$ e $\tau'$ são dois perfis distintos. Agora $N \cup {A}$ diferencia $\mathcal{O} \cup \{\tau\}$, mas pode não ser mais conjunto árvore. Isto pode acontecer caso o $A$ não seja encaixado com alguns elementos de $\mathcal{N}$. Seja $B$ um desses elementos. Vamos resolver esse não encaixamento substituindo $A$ ou $B$ por algum dos conjuntos dentre $\{A \cap B, \lnot A \cap B, A \cap \lnot B, \lnot A \cap \lnot B\}$, que são os conjuntos que chamamos de cantos de $A$ e $B$. Cantos funcionam bem pois eles não vão atrapalhar nenhum encaixamento, como vamos provar ao final. Não necessariamente todos os cantos vão estar em $S$, mas alguns pelo menos vão estar, pelo fato de $S$ ser sub-Boole, e isso será suficiente.  Para decidirmos qual substituição vamos fazer, precisamos escolher a que não vai nos fazer perder a diferenciação de nenhum perfil. Vamos por casos. Primeiramente, suponha $\tau(B) = 1$.

Suponha que $A \cap B \in S$. Como $\tau(B) = 1$, e $A$ diferencia $\tau$ e $\tau'$, então $A \cap B$ também diferencia $\tau$ de $\tau'$. Construa $N'$ substituindo $A$ por $A\cap B$ em $N \cup {A}$, isto é, $N' = N \cup {A \cap B}$. $N'$ continua a diferenciar $\mathcal{O} \cup \{\tau\}$.

Suponha que $A \cap B \notin S$. Então por $S$ ser sub-Boole, $A \cup B \in S$. A união não é suficiente para diferenciar $\tau$ e $\tau'$, então agora não podemos fazer o mesmo truque do caso anterior sem perder a diferenciação de $\tau$. E se substituissemos o $B$? Possivelmente poderiamos perder a diferenciação de algum outro par $\tilde{\tau}$ e $\tilde{\tau}'$ em $\mathcal{O}$. De fato, se existir um par destes, ele será único, como provaremos ao final em um lema. Então precisamos apenas nos preocupar com este par $\tilde{\tau}$ e $\tilde{\tau}'$. Se o $A$ diferencia este par, então podemos simplesmente definir $N' = N \setminus B \cup {A}$, que irá continuar a diferenciar $\mathcal{O} \cup \{\tau\}$ com um não encaixamento a menos. Caso contrário, $A$ pode valer ou $0$ em ambos ou $1$ em ambos. Se ele vale $0$ em ambos, então $A \cup B$ diferencia $\tilde{\tau}$ e $\tilde{\tau}'$. Defina neste caso $N' = N \setminus B \cup {A, A \cup B}$. Se ele vale $1$ em ambos, vamos ter que procurar novos cantos. Olhe para o par $A, \lnot B$. Se a interseção $A \cap \lnot B \in S$, então ela diferencia $\tilde{\tau}$ e $\tilde{\tau}'$ e podemos definir $N' = N \setminus B \cup {A, A \cap \lnot B}$. Se $A \cap \lnot B \notin S$, então a união $A \cup \lnot B \in S$. Neste caso não temos a diferenciação do $\tilde{\tau}$ e $\tilde{\tau}'$, mas temos a diferenciação do $\tau$ e do $\tau'$, então podemos voltar a ideia anterior e subtituir o $A$! Defina daí $N' = N \cup \{A \cup \lnot B\}$.

Começamos a argumentação dos dois parágrafos anteriores supondo que $\tau(B) = 1$. Se este não for o caso, troque $B$ por $\lnot B$ na argumentação inteira. 

Finalmente, em qualquer um dos casos, chegamos a um $N'$ que diferencia $\mathcal{O} \cup \{\tau\}$ e que substitui dentre dois $A,B$ não encaixados, um deles por um do seus cantos.

\begin{lema*}
  Tome dois conjuntos $A,B$ não encaixados e um $C$ encaixado com $A$ ou com $B$. Então $C$ está encaixado com qualquer um dos cantos $\{A \cap B, A \cap \lnot B, \lnot A \cap B, \lnot A \cap \lnot B\}$.
\end{lema*}

Com este lema, garantimos que nosso $N'$ tem exatamente um não encaixamento a menos que $N$. Repita este processo até não ter mais não encaixamentos, e assim chegará a um conjunto árvore que diferencia $n$ perfis, finalizando nosso passo indutivo.

Resta apenas demonstrar um fato afirmado durante a demonstração e não provado: que existe um único par $\tilde{\tau}$ e $\tilde{\tau}'$ em $\mathcal{O}$ diferenciado por $B$ e não diferenciado por nenhum outro elemento de $N$.

\begin{lema*}
  Seja $N$ um conjunto árvore, $B \in N$ e $\mathcal{O}$ um conjunto de perfis diferenciados por $N$. Se existe um par de perfis $\tilde{\tau}$ e $\tilde{\tau}'$ em $\mathcal{O}$ diferenciado apenas por $B$ e por mais ninguém, então este par é único.
\end{lema*}

\end{prova}

Usando nossa forma canônica para $S$, o teorema é equivalente a demonstrar que, dado uma família de subconjuntos $S \subset \mathcal{P}(V)$ que é sub-Boole e separa cada elemento de $V$, existe um conjunto árvore $N \subset S$ que ainda separa cada elemento de $V$.

\begin{pergunta}
  Quando $S$ é infinito, o que precisamos pedir aqui para garantir a existência de $N$?
\end{pergunta}

\begin{pergunta}
Existe uma condição topológica em $\text{Per}(S)$ que é equivalente a $N$ exibir todos os perfis? Isto ajuda a generalizar o teorema para o caso infinito?
\end{pergunta}


\begin{proposicao*}[Árvore de emaranhados não vale no caso de álgebra de Boole não-enumerável]
  Não existe conjunto árvore que exibe todos os ultrafiltros de $\mathcal{P}(\mathbb{N})$.
\end{proposicao*}
\begin{prova}
 (Pra ser escrito - ver https://mathoverflow.net/questions/508614/how-small-can-a-set-x-subset-2-omega-be-if-it-is-able-to-distinguish-any-ult)
\end{prova}
\begin{pergunta}
  Uma álgebra de Boole enumerável tem seu espaço de Stone homeomorfo ao espaço de extremidades de uma árvore. Isso indica que talvez o teorema de árvore de emaranhados vale para álgebras de Boole enumeráveis?
\end{pergunta}

\begin{pergunta}
  (Pergunta do Gustavo) Todo conjunto aberto tem display? Todo $G_{\delta}$ tem display?
\end{pergunta}

\section{Teorema da Dualidade Árvore-Emaranhado}

Seja $\mathcal{B}$ uma álgebra de Boole e considere $\mathcal{E} \subset \mathcal{P}(\mathcal{B})$.Chamamos os elementos de $\mathcal{E}$ de estrelas. Dizemos que duas estrelas $E_1,E_2 \in \mathcal{E}$ se colam se existe $a \in E_1, \lnot a \in E_2$.

Vamos brincar aqui de colar estas estrelas com o objetivo de que todo $A$ cole com algum $\lnot A$ eventualmente. Vamos chamar de solução um subconjunto $T \subset \mathcal{E}$ tal que $\cup T$ é fechado pra complementar.

Vamos dar aqui caracterizações para a existência de soluções para um dado $\mathcal{E}$, usando orientações e emaranhados. Lembrando, dado algum $S \subset \mathcal{B}$, uma orientação $\tau$ de $S$ é um mapa $\tau: S \rightarrow \{0,1\}$ tal que $\tau(A) = 0 \iff \tau(\lnot A) = 1$. Um emaranhado é uma orientação $\tau$ tal que $A\subset B$ implica $\tau(A) \leq \tau(B)$ - colocado de outra forma, estamos dizendo que se $\tau(A) = 1$ e $A \subset B$ então $\tau(B) = 1$. Vamos nos referir a esta propriedade como consistência - um emaranhado é, então, uma orientação consistente. Dizemos que uma orientação ou emaranhado $\tau$ satisfaz um $T \subset S$ se $\tau(t) = 1$ pra todo $t \in T$.

\begin{proposicao}[Ida da Dualidade Fraca]
  Seja $\mathcal{B}$ uma álgebra de Boole e considere $\mathcal{E} \subset \mathcal{P}(\mathcal{B})$. Seja $S$ o fecho pra complementar de $\cup \mathcal{E}$. Suponha que $S$ seja finito. Se $\mathcal{E}$ não admite solução, isto é, se não existe $\mathcal{T} \subset \mathcal{E}$ tal que $\cup \mathcal{T}$ é fechado pra complementar, então existe uma orientação $\tau$ de $S$ tal que pra todo $E \in \mathcal{E}$, temos que $\tau$ não satisfaz $E$.
\end{proposicao}
\begin{prova}
Suponha que $\mathcal{E}$ não admite solução.

Dado um conjunto de estrelas $\mathcal{F} \subset \mathcal{E}$, vamos colar todas estas estrelas e coletar em um conjunto os elementos delas que não foram pareados com seu complementar. Precisamente, definimos $ \mathcal{F}^* \defeq \{a \in \cup \mathcal{F} \mid \lnot a \notin \cup \mathcal{F}\}$. Agora vamos extender nosso conjunto de estrelas $\mathcal{E}$ para um maior $\overline{\mathcal{E}}$ com todos estes $\mathcal{F}^*$ - pensamos nisso como um fecho para a operação de colagem (veja que esta operação, como definimos, pode colar infinitas estrelas de uma vez). Isto é, $\overline{\mathcal{E}} \defeq \mathcal{E} \cup \{\mathcal{F}^* \mid \mathcal{F} \subset \mathcal{E}\}$.  Note que $\varnothing \in \overline{\mathcal{E}}$ se e somente se existe um $\mathcal{T} \subset \mathcal{E}$ tal que $\cup \mathcal{T}$ é fechado pra complementar. 

Para um dado $X \subset S$, considere a seguinte propriedade

\[
(\star) \text{ para cada } E \in \overline{\mathcal{E}}\text{ existe } a \in E \text{ tal que } \lnot a \in X
\]

Defina $P = \{a \in S \mid \{a\} \notin \overline{\mathcal{E}}\}$. Vamos mostrar que $P$ satisfaz $(\star)$. Suponha que não - isto é, suponha que exista um $E \in \overline{\mathcal{E}}$ tal que pra todo $a \in E$, temos $\lnot a \notin P$, isto é, $\{\lnot a\} \in \overline{\mathcal{E}}\ $. Se $\{\lnot a\} \in \mathcal{E}$, defina $\mathcal{F}_a = \{\{\lnot a\}\}$ e se $\{\lnot a\} = \mathcal{F}^*$, defina $\mathcal{F}_a = \mathcal{F}$. Com isso, defina $\mathcal{T} \defeq \{E\} \cup \underset{a \in E}{\bigcup} \mathcal{F}_a$ quando $E \in \mathcal{E}$ e $\mathcal{T} \defeq \mathcal{F}_E \cup \underset{a \in E}{\bigcup} \mathcal{F}_a$ quando $E = \mathcal{F}_E^*$. Em ambos os casos, $\mathcal{T}$ é tal que $\cup \mathcal{T}$ é fechado pra complementar. Como supomos que $\mathcal{E}$ não admite solução, então temos uma contradição e podemos concluir que $(\star)$ vale para $P$.

Agora, suponha que $a,\lnot a \in P$. Então não pode ser o caso que ambos $P \setminus {a}$ e $P \setminus {\lnot a}$ não satisfazem $(\star)$. Pois se isto acontecesse, teriamos duas estrelas $E_1, E_2 \in \overline{\mathcal{E}}$ tais que pra todo $b \in E_1$ temos $\lnot b \notin P \setminus {a}$ e pra todo $c \in E_2$ temos $\lnot c \notin P \setminus {a}$. Mas como $P$ satisfaz $(\star)$, existe um $x \in E_1$ tal que $\lnot x \in P$, e como $\lnot x \notin P \setminus {a}$, vamos ter que $\lnot x = a$, isto é, $\lnot a \in E_1$. Analogamente, temos que $a \in E_2$. Assim, $\{E_1,E_2\}^*$ seria uma contradição para propriedade $(\star)$ de $P$.

Podemos então sucessivamente ir retirando um dos $a,\lnot a$ de $P$, em cada passo mantendo a propriedade $(\star)$ e eventualmente chegar em um conjunto $O \subset P$ que não contém dois elementos inversos um do outro. Colocado de outra forma, $O$ é o conjunto minimal em relação a inclusão que satisfaz $(\star)$ dentre os subconjuntos de $P$. A finitude de $S$ garante que este minimal existe. A prova está finalizada aqui: a orientação $\tau$ é a que valora todos os elementos de $O$ em $1$, seus inversos em $0$, e escolhe arbitrariamente entre $a$ e $\lnot a$ para elementos que não estão em $O$ nem tem seus inversos em $O$. A propriedade $(\star)$ de $O$ garante que $\tau$ não vai satisfazer nenhuma estrela de $\mathcal{E}$.

\end{prova}

Queremos agora trocar `orientação' por `emaranhado' neste último resultado. É com essa troca que vamos conseguir aplicações em grafos, como os teorema de dualidade entre decomposição de árvore e bramble (e possivelmente mais aplicações em lógica/conjuntos?). Pra isso, vamos ter que fortalecer nossas hipóteses, pedindo que o conjunto de estrelas $\mathcal{E}$ tenha uma certa propriedade, modelada em algo que vale para o caso de separações de grafos. Esta propriedade vai garantir que conseguiremos adaptar nossa construção da demonstração anterior de forma que a orientação se torne consistente, isto é, defina um emaranhado.

A propriedade é a seguinte. Diga que $\mathcal{E}$ satisfaz a hipótese da dualidade se para quaisquer duas estrelas $E,F \in \mathcal{E}$, com $E = \{a, a_1, \dots, a_n\}, F = \{b, b_1, \dots, b_m\}$ satisfazendo $\lnot a \subset b$, vamos ter que existe um $x \in S$ tal que $\lnot b \subset x \subset a$ e existem estrelas $E',F' \in \mathcal{E}$ com $E' = \{x, a_1',\dots, a_n'\}, F' = \{\lnot x, b_1', \dots, b_m'\}$ e $a_i \subset a_i', b_i \subset b_i'$.

\begin{proposicao}[Ida da Dualidade Forte]
  Seja $\mathcal{B}$ uma álgebra de Boole e considere $\mathcal{E} \subset \mathcal{P}(\mathcal{B})$. Seja $S$ o fecho pra complementar de $\cup \mathcal{E}$. Suponha que $S$ seja finito. Suponha que $\mathcal{E}$ satisfaz a hipótese da dualidade. Se $\mathcal{E}$ não admite solução, isto é, se não existe $\mathcal{T} \subset \mathcal{E}$ tal que $\cup \mathcal{T}$ é fechado pra complementar, então existe um emaranhado $\tau$ de $S$ tal que pra todo $E \in \mathcal{E}$, temos que $\tau$ não satisfaz $E$.
\end{proposicao}

\begin{prova}
  Suponha que $\mathcal{E}$ não admita solução. Vamos traçar os mesmos passos da demonstração da ida da dualidade fraca, definindo $\overline{\mathcal{E}}$, $P$, e a propriedade $(\star)$, que $P$ irá satisfazer. Agora, para construirmos nosso $O \subset P$ que irá dar origem ao emaranhado, precisamos tomar mais cuidado para garantir que ele será consistente. Para isso, defina uma propriedade a mais: para um $X \subset S$ qualquer, considere

  \[
  (\bullet) \text{ se } a \in X \text{ e } a \subset b \text{ para um } b \in P \text{ então } b \in X
  \]

  Claramente $P$ satisfaz $\bullet$. A finitude de $S$ permite que tomemos um $O \subset P$ que é minimal em $\{X \subset P \mid X \text{ satisfaz } (\star) \text{ e } (\bullet)\}$.

  Vamos mostrar que não pode existir em $O$ dois elementos $r,s$ tais que $r \subset \lnot s$. Isto vai nos permitir concluir, ao mesmo tempo, que $O$ vai induzir uma orientação (pois não pode existir $\{a,\lnot a\} \subset O$) e que essa orientação vai ser consistente.

  Tome, então, $r,s \in O$ com $r \subset \lnot s$. Se existisse $r',s' \in O$ com $r' \subset r$ e $s' \subset s$, então $r' \subset \lnot s'$ - usando a finitude de $S$ novamente, podemos então supor que $r,s$ foram tomados minimais. 

  $O \setminus r$ não pode satisfazer ambos $(\star)$ e $(\bullet)$ devido a minimalidade de $O$. Se $O \setminus r$ não satisfizesse $(\bullet)$, deveria existir um $a \in O\setminus r$ com $a \subset b$, $b \in P$ e $b \notin O\setminus r$. Como $a \in O$ e $O$ satisfaz $(\bullet)$, temos que $b \in O$ e portanto $b = r$. Como $r$ é minimal, então $a = r$, contradizendo $\ \in O \setminus r$. Concluimos que $O \setminus r$ não satisfaz $(\star)$, e igualmente também $O \setminus s$ não satisfaz $(\star)$.

  Tome então $E_r,E_s \in \overline{\mathcal{E}}$ tais que pra todo $a \in E_r$, temos $\lnot a \notin O \setminus r$ e todo $b \in E_s$, temos $\lnot b \notin O \setminus s$. Como $O$ satisfaz $(\star)$, deve existir um $x \in E_r,y \in E_s$ tais que $\lnot x \in O, \lnot y \in O$ - e portanto $x = \lnot r, y = \lnot s$. Chame $E_r = \{\lnot r, a_1,\dots, a_n\}$ e $E_s = \{\lnot s, b_1, \dots, b_m\}$. Se algum dos $\lnot a_i$ não estivesse em $P$, então isso significaria que $\{\lnot a_i\} \in \overline{\mathcal{E}}$, e aí poderíamos fazer a colagem de $E_r$ com $\{\lnot a_i\}$ para obter que $E_r \setminus \{a_i\} \in \overline{\mathcal{E}}$ - e este ainda é uma testemunha para $O \setminus r$ não satisfazer $(\star)$. Podemos supor, então, que realizamos este processo a quantidade de vezes necessária para chegarmos em um $E_r$ tal que seus elementos $a_i$ são todos tais que $\lnot a_i \in P$. Igualmente, tomamos o $E_s$ tal que todos os $b_i$  estão em $P$.

  Finalmente vamos utilizar a hipótese de dualidade. Porém, precisamos utilizá-la referente ao $\overline{E}$. Mais tarde voltamos a este detalhe - suponha por agora que $\overline{E}$ satisfaz a hipótese de dualidade.

 Lembrando que $r \subset \lnot s$, aplique a hipótese de dualidade em $E_r$ e $E_s$. Obtemos assim, duas estrelas $E_r' = \{a,a_1',\dots,a_n'\}, E_s' = \{\lnot a, b_1', \dots, b_m'\}$ em $\overline{\mathcal{E}}$ tais que $\lnot s \subset a \subset \lnot r$, $a_i \subset a_i'$, $b_i \subset b_i'$. Veja que $\lnot a_i' \notin O$, pois se esse não fosse o caso, então $\lnot a_i$ seria um elemento de $P$ que contém $\lnot a_i \in O$, e pela propriedade $(\bullet)$, teriamos $\lnot a_i \in O$, que já sabemos que não vale. Igualmente, $\lnot b_i' \not in O$. Como a colagem de $E_r'$ com $E_s'$ resulta em uma estrela de $\overline{\mathcal{E}}$ cujos elementos estão dentre os $a_i',b_i'$, ela seria então uma testemunha para $O$ não satisfazer $(\star)$ - contradição.

  Assim, chegamos a conclusão que $O$ é um subconjunto de $S$ tal que $a \in O \implies \lnot a \notin O$ e que se fecharmos $O$ para cima em relação a inclusão, isto é, $\overline{O} = \{x \in S \mid a \subset x \text{ para algum } a \in O\}$, então este fecho ainda satisfaz $x \in \overline{O} \implies \lnot x \notin \overline{O}$. Definimos $\tau$ como $1$ para todo elemento de $\overline{O}$, e $0$ para todos os seus inversos. Para o resto de $S$, extenda $\tau$ de qualquer maneira que o mantenha consistente. Dessa forma, chegamos a um emaranhado $\tau$ que devido a propriedade $(\star)$ de $O$, não irá satisfazer nenhuma estrela de $\mathcal{E}$.

  Voltemos a ponta solta de alguns parágrafos atrás, e vamos mostrar que $\mathcal{E}$ satisfazer a hipótese da dualidade implica que $\overline{\mathcal{E}}$ a satisfaz também.

  ...

\end{prova}

Mostramos que a não existência de uma orientação que não satisfaz nenhuma estrela de $\mathcal{E}$ é uma condição suficiente para garantirmos a existência de uma solução. Mas não é uma condição necessária: ciclos dão origem a soluções, e também a orientações que não satisfazem estrelas. Por exemplo, se tomarmos $\mathcal{E} = \{\{a,\lnot b \}, \{b, \lnot c\}, \{c, \lnot a\}\}$, então a união de todas estas estrelas é fechado para complementar, e a orientação $\tau$ que escolhe $\tau(b) = 1, \tau (c) = 1, \tau (a) = 1$ não satisfaz nenhuma estrela.

Formalmente, definimos um ciclo em $\mathcal{E}$ como sendo uma sequência de estrelas $E_1, \dots, E_n$ em $\mathcal{E}$ tal que $E_i$ cola com $E_{i+1}$ e $E_n$ cola com $E_1$.

\begin{proposicao}[Volta da Dualidade Fraca e Forte]
  Seja $\mathcal{B}$ uma álgebra de Boole e considere $\mathcal{E} \subset \mathcal{P}(\mathcal{B})$. Seja $S$ o fecho pra complementar de $\cup \mathcal{E}$. Suponha que $S$ seja finito. Se $\mathcal{E}$ admite solução, isto é, se existe um $\mathcal{T} \subset \mathcal{E}$ tal que $\cup \mathcal{T}$ é fechado pra complementar, então ou existe um ciclo em $\mathcal{E}$ ou toda orientação $\tau$ (e portanto todo emaranhado) de $S$ satisfaz algum $E \in \mathcal{E}$.
\end{proposicao}
\begin{prova}
Suponha que temos uma solução $\mathcal{T}$ e que $\mathcal{E}$ não admita ciclos. Tome uma orientação $\tau$ de $S$ qualquer.

...
\end{prova}

Juntando a ida e a volta, obtemos os resultados que queremos.

\begin{teorema*}[Dualidade Árvore-Emaranhado Fraca]
  Seja $\mathcal{B}$ uma álgebra de Boole e considere $\mathcal{E} \subset \mathcal{P}(\mathcal{B})$. Seja $S$ o fecho pra complementar de $\cup \mathcal{E}$. Suponha que $S$ seja finito, e que $\mathcal{E}$ não admita ciclos. Então vale um dos dois, e não ambos:
  
  \begin{itemize}[label=\ding{212}]
  \item existe um subconjunto $T \subset \mathcal{E}$ tal que $\cup T$ é fechado pra complementar
  \item existe uma orientação $\tau$ de $S$ que não satisfaz $E$ para toda estrela $E \in \mathcal{E}$.
\end{itemize}
\end{teorema*}
\begin{prova}
  A volta da Dualidade Forte e Fraca garante que ambos não valem, e Ida da Dualidade Fraca garante que pelo menos um vale.
\end{prova}

\begin{teorema*}[Dualidade Árvore-Emaranhado Forte]
  Seja $\mathcal{B}$ uma álgebra de Boole e considere $\mathcal{E} \subset \mathcal{P}(\mathcal{B})$. Seja $S$ o fecho pra complementar de $\cup \mathcal{E}$. Suponha que $S$ seja finito, e que $\mathcal{E}$ não admita ciclos. Suponha que $\mathcal{E}$ satisfaz a hipótese da dualidade. Então vale um dos dois, e não ambos:
  
  \begin{itemize}[label=\ding{212}]
  \item existe um subconjunto $T \subset \mathcal{E}$ tal que $\cup T$ é fechado pra complementar
  \item existe um emaranhado $\tau$ de $S$ que não satisfaz $E$ para toda estrela $E \in \mathcal{E}$.
\end{itemize}
\end{teorema*}
\begin{prova}
  A volta da Dualidade Forte e Fraca garante que ambos não valem, e Ida da Dualidade Forte garante que pelo menos um vale.
\end{prova}


%Fazer caso forte

%Mostrar que minha versão implica a versão do Diestel - ou, melhor talvez, mostrar que implica no tree-decomp e bramble tree width duality?



%Generalização pro infinito? (arjuna ?)

%Ideia do aurichi de pensar classes de grafos caracterizadas por conjuntos de estrelas? (todos os grafos que nascem colando estrelas desse conjunto)

%Escrever versão bonitinha como puzzle



  \section{A Lógica dos Emaranhados}

  Álgebras de Boole são o estudo de teorias do cálculo proposiconal. Vou escrever sobre esta equivalência mais pra frente - mas um resumo: considere um alfabeto $\Sigma = \{p_1,p_2, \dots\}$ que vamos tratar como as proposições atômicas. $\Sigma_n$ vai ser o alfabeto com $n$ proposições para um $n \in w$ e $\Sigma_w$ o alfabeto com enumeráveis proposições. Chame de $\text{Sent}(\Sigma)$ o conjunto de sentenças, como usualmente definido, do cálculo proposicional. Dado um conjunto $T \subset \text{Sent}(\Sigma)$, podemos definir a relação de equivalência $\phi \cong_T \psi$ quando $T \vdash \phi \leftrightarrow \psi$. O quociente de $\text{Sent}(\Sigma)$ por $\cong_T$ vai ser chamado de $\mathcal{B}_T$ e os operadores $\vee,\wedge, \lnot$ o mune de uma estrutura de álgebra de Boole. Quando $T = \varnothing$, $\mathcal{B}_T$ é livre sobre $\Sigma$. Variando $T$, obtemos todas as álgebras de Boole (veja que isto é equivalente a dizer que toda álgebra de Boole é quociente de uma livre, isto é, achar uma apresentação pra uma álgebra de Boole). Ultrafiltros de $\mathcal{B}_T$ correspondem a modelos da teoria $T$, isto é, valorações (consistentes) das proposições que valoram $T$ em $1$.

Dado um $S \subset \mathcal{B}_T$, um perfil $\tau$ de $S$ é uma valoração das proposições de $S$ que satisfazem consistência apenas em relação aos elementos que $S$ contém. Por exemplo, $\tau$ pode valorar três proposições $A,B,C \in S$ com $1$, mesmo que $A \cap B \cap C = \varnothing$, desde que em $S$ não existam as interseções $A \cap B$, $B \cap C$ e $A \cap C$. Interpretamos intuitivamente que $\tau$ é `$S$-consistente', e portanto não necessariamente pode ser extendido a um modelo da teoria inteira. Podemos ver perfis, então, como uma forma de generalização de valorações - uma nova semântica pro cálculo proposicional.

Dois resultados clássicos da teoria de cálculo proposicional relacionam a semântica do cálculo (modelos, valorações) com a sintaxe (dedução natural, ou sistema axiomático com regras de inferência): o teorema da completude e o teorema da compacidade. Estes teoremas valem pra perfis? 

Vou escrever definições para álgebras de Boole e que acredito que são naturalmente traduzidas para a linguagem de cálculo proposicional.


\newcommand{\Sproves}[1]{\underset{#1}{\vdash}}

\begin{definicao*}[$S$-provas, $S$-possibilidade e $S$-consistência]


  Seja $\mathcal{B}$ uma álgebra de Boole, $S \subset \mathcal{B}$ um subconjunto fechado para complementar com $\varnothing \in S$ e $T \subset S$ um subconjunto qualquer. \\
  
  Dizemos que $T$ $S$-prova um $A \in S$ e escrevemos $T \Sproves{S} A$ se vale pelo menos um dos itens:
    
  \begin{itemize}[label=\ding{212}]
   \item $A \in T$
    \item (conjunção) $A = B \cap C$ para $B,C \in S$ onde $T \Sproves{S} B$ e $T \Sproves{S} C$
    \item (implicação) existe um $B \in S$ tal que $T \Sproves{S} B$ e vale que ou $B \subset A$ ou $T \Sproves{S} \lnot B \cup A$
    \item (dedução) $T \cup \{\lnot A\} \Sproves{S} \varnothing $
  \end{itemize} 

  Quando não houver ambiguidade sobre quem é o $S$ subjacente, denotamos $T_{\vdash}$ como o conjunto de $A \in S$ tais que $T \Sproves{S} A$.\\

  Dizemos que $T$ é $S$-consistente se não vale que $T \Sproves{S} \varnothing$, ou seja, $\varnothing \notin T_{\vdash}$.
\end{definicao*}

Esta definição foi inspirada na definição de dedução natural em cálculo proposicional. Quando $T \Sproves{S} A$ podemos desenhar uma árvore finita que explicita a estrutura recursiva dos usos de conjunção, implicação e dedução. Em resumo, $T$ $S$-prova $A$ quando existe uma dedução natural de $A$ a partir de $T$ que usa apenas elementos de $S$.

\begin{pergunta}
  Originalmente eu tinha definido uma noção mais fraca de $S$-prova que não era suficiente para equivalência sintáxe-semântica. Essa noção era mais inspirada na definição de sistema axiomático do cálculo proposocional.Esses outros sistemas axiomáticos clássicos que são equivalentes a dedução natural - se definirmos uma $S$-prova como sendo uma prova (a partir destes outros sistemas) que usa apenas elementos de $S$, isso vai ser equivalente a definição que acabamos de dar? Penso que não, por que para se provar equivalência acho que vamos preicsar usar conjunções e implicações que não estão em $S$.
\end{pergunta}

É preciso mostrar uma certa noção de `boa fundação' para $\Sproves{S}$ para justificar nosso uso de indução na próxima proposição.

\begin{proposicao*}($S$-corretude)
  Seja $\mathcal{B}$ uma álgebra de Boole, $S \subset \mathcal{B}$ um subconjunto fechado para complementar com $\varnothing \in S$ e $T \subset S$ um subconjunto qualquer. \\

  Se existe um perfil $\tau$ de $S$ que valora todos os elementos de $T$ em $1$, então $T$ é $S$-consistente.
\end{proposicao*}

\begin{prova}
  Tome um perfil $\tau$ de $S$ que valora todos os elementos de $T$ em $1$. Por indução podemos deduzir que $\tau(A) = 1$ para todo $A \in S$ tal que $T \Sproves{S} A$. Como $\tau(\varnothing) = 0$, não podemos ter $T \Sproves{S} \varnothing$.
\end{prova}

% \begin{exemplo*}[Um $S$ que não satisfaz dedução nem completude]
% Seja $\mathcal{B}$ o conjunto de partes de $V = \{1,2,3,4\}$ e
% \begin{align*}
% X &= \{1, 2\}           &  \lnot X &= \{3,4\}              \\
% Y &= \{2, 4 \}         &  \lnot Y &= \{1,3\}  \\
% Z&= \{1\}          &  \lnot Z &= \{2,3,4\} \\
% W &= \{4\}         &   \lnot W &= \{1,2,3\} \\
% A &= \{1,3\}       &    \lnot A &= \{2,4\}    
% \end{align*}


% Colete todos estes conjuntos no $S = \{X,Y,Z,W, \lnot X, \lnot Y, \lnot Z, \lnot W, \lnot A, \varnothing, V \}$, e considere $T = \{X,Y\}$. Seja $T_{\vdash}$ o conjunto de elementos $S$-provados por $T$. É uma continha ver que $T_{\vdash} = \{X,Y, \lnot Z, \lnot w\}$. Logo $T$ é $S$-consistente.

% Agora, $T \cup \{A\} \Sproves{S} Z$ pois $Z = X \cap A$. Daí $Z \cap Y = \varnothing$ e logo $T \cup \{A\} \Sproves{S} \varnothing$.

% Similarmente, $T \cup \{\lnot A\} \Sproves{S} W$ pois $W = \lnot A \cap Y$. Daí $W \cap X = \varnothing$ e logo $T \cup \{\lnot A\} \Sproves {S} \varnothing$.

% Assim, $T$ não satisfaz o teorema da dedução, pois temos ambos $T \cup \{\lnot A\} \Sproves {S} \varnothing$ e $T \cup \{A\} \Sproves{S} \varnothing$ mas $A,\lnot A \notin T_\vdash$. E também não satisfaz o teorema da completude: não pode haver perfil $\tau$ valorando $T$ inteiramente em $1$, pois $\tau$ valoraria um dos $A, \lnot A$ como $1$ e aí por $S$-corretude não poderiamos ter $T \cup \{A\}$ e $T \cup \{\lnot A\}$ ambos inconsistentes.

% \end{exemplo*}

% \begin{definicao*}[$S$-dedução]
%        Seja $\mathcal{B}$ uma álgebra de Boole, $S \subset \mathcal{B}$ um subconjunto fechado para complementar com $\varnothing \in S$.

%        Suponha que pra qualquer $T \subset S$, temos que se $T \cup \{A\} \Sproves{S} \varnothing$ então $T \Sproves{S} \lnot A$. Dizemos neste caso que $S$ admite $S$-dedução.
% \end{definicao*}

\begin{proposicao*}[$S$-completude]
   Seja $\mathcal{B}$ uma álgebra de Boole, $S \subset \mathcal{B}$ um subconjunto fechado para complementar com $\varnothing \in S$ e $T \subset S$ um subconjunto qualquer. \\
   
   Se $T$ é $S$-consistente, deve existir um perfil $\tau$ de $S$ que valora todos os elementos de $T$ em $1$.
\end{proposicao*}
\begin{prova}
  Suponha que $T$ é $S$-consistente. Seja $\tilde{T} = \{A \mid T \Sproves{S} A\}$. Veja que $\tilde{T} \Sproves{S} B$ se e somente se $T \Sproves{S} B$ se e somente se $B \in \tilde{T}$ - isto é, $\tilde{T}$ é fechado para $S$-provabilidade. A $S$-consistência de $T$ garante que $\varnothing \notin \tilde{T}$ - o que garante também que se $A \in \tilde{T}$, então $\lnot A \notin \tilde{T}$.\\
   
  Suponha que exista um $B \in S$ tal que $B, \lnot B \notin \tilde{T}$. Devemos ter que $\tilde{T} \cup \{B\}$ ainda é $S$-consistente, pois se não fosse, por dedução teriamos que $\tilde{T} \Sproves{S} \lnot B$. Extendemos $\tilde{T} \cup \{B\}$ para um $\widetilde{\tilde{T} \cup \{B\}}$ fechado para $S$-provabilidade, que não contém o vazio. Se ainda existir um $C \in S$ tal que nem $C$ nem $\lnot C$ estão em $\widetilde{\tilde{T} \cup \{B\}}$, repetimos o processo. Assim, por indução transfinita (ou lema de Zorn), chegamos a um elemento $U \subset S$ que: contém $T$, é $S$-consistente, é fechado pra $S$-provabilidade, e é maximal em relação a estas três propriedades. A maximalidade de $U$ garante que ele deve conter para todos $A \in S$ exatamente um dentre $A, \lnot A$. Definindo $\tau$ como sendo $1$ para todos os elementos de $U$ e $0$ fora, temos nosso perfil desejado.
\end{prova}

\newcommand{\Simplies}[1]{\underset{#1}{\vDash}}

\begin{definicao*}[Satisfação, implicação semântica]
Seja $\mathcal{B}$ uma álgebra de Boole, $S \subset \mathcal{B}$ um subconjunto fechado para complementar com $\varnothing \in S$ e $T \subset S$ um subconjunto qualquer.\\

Se um perfil $\tau$ de $S$ valora todo elemento de $T$ como $1$, dizemos que $\tau$ satisfaz $T$.\\


Dado $A \in S$, dizemos que $T$ implica semanticamente $A$, ou $A$ é consequência semântica de $T$, e escrevemos $T \Simplies{S} A$, se pra todo perfil $\tau$ de $S$ que satisfaz $T$, temos que $\tau(A) = 1$.
\end{definicao*}

A equivalência síntaxe-semântica é a junção de corretude e completude.

\begin{proposicao*}[$S$-equivalência sintáxe-semântica]
  Seja $\mathcal{B}$ uma álgebra de Boole, $S \subset \mathcal{B}$ um subconjunto fechado para complementar com $\varnothing \in S$ e $T \subset S$ um subconjunto qualquer. Então, pra todo $A \in S$,

       \[T \Sproves{S} A \iff T \Simplies{S} A \]
\end{proposicao*}
\begin{prova}
($\implies$) Na demonstração de corretude, argumentamos por indução que todo perfil que satisfaz $T$ deve satisfazer $A$. Também podemos mostrar o mesmo resultado usando diretamente o enunciado de coretude: suponha que $T \Sproves{S}$. Se existisse um perfil $\tau$ que satisfaz $T$ mas não satisfaz $A$, então $\tau$ satisfaria $T \cup \{\lnot A\}$. Por $S$-corretude, $T \cup \{\lnot A\}$ deveria ser consistente, mas $T \Sproves{S} A$ implica que $T \cup \{\lnot A\}$ é inconsistente.

($\impliedby$) Suponha que não vale que $T \Sproves{S} A$. Então $T \cup \{\lnot A\}$ é consistente, pois se não fosse teriamos por dedução que $T \Sproves{S} A$. Por $S$-completude, deve existir um perfil que satisfaz $T$ mas não satisfaz $A$, logo não vale que $T \Simplies{S} A$.
\end{prova}


% \begin{exemplo*}[Um $S$ satisfazendo completude e dedução, mas não equivalência síntaxe-semântica]
%   Seja $\mathcal{B}$ o conjunto de partes de $V = \{a,b,1,2,3,4\}$ e
% \begin{align*}
% A &= \{a, 1, 2\}           &  \lnot A&= \{3,4,b\}              \\
% B &= \{a, 3, 4 \}         &  \lnot B &= \{1,2,b\}  \\
% C&= \{1,2\}          &  \lnot C &= \{a,b,3,4\} 
% \end{align*}

% Faça $S = \{A,B,C, \lnot A, \lnot B, \lnot C, \varnothing, V\}$ e $T = \{A,\lnot C\}$. Veja que $T_{\vdash} = T$, pois $A,C$ não possuem interseção em $S$ e não estão contidos em ninguém. Como $C = A \cap \lnot B$, qualquer perfil $\tau$ que valore $A, \lnot B$ como $1$ deve valorar $C$ como $1$, logo $T \Simplies{S} B$. Assim, não temos equivalência sintática-semântica.\\

% Também conseguimos mostrar que $S$ satisfaz $S$-dedução. Suponha que $T \cup \{X\} \Sproves{S} \varnothing$. Se $T$ é inconsistente, então $T \Sproves{S} \lnot X$ pois $T$ prova qualquer coisa ($\varnothing$ está contido em qualquer conjunto). Se $T$ não é inconsistente, então há alguns casos a considerar. Se $X = \lnot Y$ para algum $Y \in T$, então é claro $T \Sproves{S} \lnot X$. Fora as interseções entre um conjunto e seu complementar, temos apenas duas outras interseções resultando no vazio: $C \cap B = C \cap \lnot A = \varnothing$.Portanto, $X$ deve ser um dentre $C, B, \lnot A$. Se $X = C$, devemos ter ou $B$ ou $\lnot A$ em $T$ - ambos estão contidos em $\lnot C$ e logo $T \Sproves{S} \lnot C$. Se $X = B$ ou $X = \lnot A$, devemos ter $C \in T$, e $C$ está contido em $A$ e $\lnot B$, logo $T \Sproves{S} \lnot X$ em ambos os casos.\\
% \begin{teorema*}[$S$-completude]
%   Seja $\mathcal{B}$ uma álgebra de Boole, $S \subset \mathcal{B}$ um subconjunto fechado para complementar com $\varnothing \in S$ e $T \subset S$ um subconjunto qualquer. 

%   Então, existe um perfil $\tau$ de $S$ que valora todos os elementos de $T$ em $1$ se e somente se $T$ é $S$-consistente.
% \end{teorema*}
% \begin{prova}
%   ($\implies$) Tome um perfil $\tau$ de $S$ que valora todos os elementos de $T$ em $1$. Por indução deduzimos que $\tau(A) = 1$ para todo $A \in S$ tal que $T \Sproves{S} A$. Se existisse $R \subset S$ tal que $T \cup R \Sproves{S} \varnothing$, então $\tau$ tem que valorar pelo menos um dos $A \in T$ como $0$ - logo não podemos ter $T \cup \{\lnot A\} \Sproves{S} \varnothing$ para todo $A \in R$.

%   ($\impliedby$) Suponha que $T$ é $S$-consistente. Seja $\tilde{T} = \{A \mid T \Sproves{S} A\}$. Veja que $\tilde{T} \Sproves{S} B$ se e somente se $T \Sproves{S} B$ se e somente se $B \in \tilde{T}$ - isto é, $\tilde{T}$ é fechado para $S$-provabilidade. A $S$-consistência de $T$ garante que $\varnothing \notin \tilde{T}$ - o que garante também que se $A \in \tilde{T}$, então $\lnot A \notin \tilde{T}$. Suponha que exista um $B \in S$ tal que $B, \lnot B \notin \tilde{T}$. Queremos mostrar que $\tilde{T} \cup \{B\}$ ainda é $S$-consistente. Tome um conjunto finito $R \subset S$ tal que $T \cup \{B\} \cup R$ consegue $S$-provar o vazio.

% \end{prova}
% Este também é um exemplo de um $S$ satisfazendo $S$-dedução que não é uma álgebra de Boole, nem uma álgebra sub-Boole.

% \end{exemplo*}


\begin{proposicao*}[$S$-compacidade]
Seja $\mathcal{B}$ uma álgebra de Boole, $S \subset \mathcal{B}$ um subconjunto fechado para complementar com $\varnothing \in S$ e $T \subset S$ um subconjunto qualquer.\\

Se pra todo subconjunto finito de $T$ existe um perfil que o satisfaz, então existe um perfil que satisfaz $T$.

\end{proposicao*}

\begin{prova}
  Se não existisse um perfil que satisfaz $T$, então por $S$-completude, $T$ deve ser inconsistente, isto é, $T \Sproves{S} \varnothing$. Deve existir uma árvore finita que atesta esta prova. O subconjunto $T'$ finito de elementos de $T$ usados nessa prova é, então, inconsistente, e por $S$-corretude, não pode ser satisfazível por nenhum perfil.
\end{prova}

\begin{pergunta}
Exemplos legais?
\end{pergunta}


Agora, vamos para lógica de primeira ordem. Seja $\mathcal{B}$ o conjunto de sentenças de uma lógica de primeira ordem quocientado pela relação de provabilidade a partir de uma teoria.



% se S tem todas as interseçoes possiveis de um conj finito de interseção vazia, então todo perfil de S tem ultrafiltro

% mostrar aquela equiv de proposicao = conj de modelos que satisfazem

% se um perfil não satisfaz uma teoria inconsistente, então ele se extende a um ultrafiltro

% S = n-cláusulas, tangles de S não tem 2n-contradições

% como tangles se relacionam com complexidade computacional? 













  %Fazer a equivalência de Lindenbaum-Tarski pra álgebras de Boole. Traduzir emaranhados para esse contexto, onde vão virar `modelos $S$-consistentes'. Generalizar essa definição para outras lógicas (não precisa passar por Lindenbaum-Tarski?). Se perguntar o que vale da teoria usal agora para essas $S$-consistentes. (compacidade, completude, etc).
  
  %Pensar o que o teorema de árvore de emaranhados e dualidade árvore-emaranhado dizem no caso de lógica, e se tem como simplificar a prova, ou generalizar.
  
  % Pensar o caso infinito. 

  % Agora ainda mais especulativo: coloca uma relação no conjunto de modelos de uma teoria (por exemplo, a relação de ser extensão de forcing) - existem tangles nesse grafo? Brisar em lógica modal e teoria do JDH de multiverso.


\bibliographystyle{plain}
\bibliography{references}


\end{document}
